% Problem record op_0af3e22238a42d74. This TeX file is derived from
% database/problems_json/op_0af3e22238a42d74.json by scripts/sync-tex.mjs; edit the JSON record
% and rerun the script rather than editing this file.
\section{Optimal dimension of a 9-qubit code that detects 1 error}
\label{sec:op_0af3e22238a42d74}

\paragraph{Problem.}
What is the largest dimension of a nine-qubit code that detects every single-qubit error exactly?

Let $P$ be the orthogonal projector onto a subspace of $(\mathbb C^2)^{\otimes9}$. Exact detection here requires $PEP=c_E P$ for a scalar $c_E$ for each $E$ that applies one Pauli operator $X$, $Y$, or $Z$ to one qubit and the identity to the others. By linearity these conditions cover arbitrary one-qubit errors. Determine the largest rank of $P$ subject to all the conditions, allowing any subspace.

\subsection*{Status}

Unsolved

\subsection*{Source}

Anglès Munné and Huber record the gap for nine qubits and distance two in Sec.~4.2, Table~1 \sourcecite{ref:d6-04-amh26}{AMH26}.

\subsection*{Progress}

\begin{itemize}
  \item Rigby, Olivier, and Jarvis give a nine-qubit code of dimension $100$ in their search for codeword-stabilized codes \sourcecite{ref:d6-04-roj19}{ROJ19}.

  \item Table~1 of Anglès Munné and Huber lists the linear programming bound $112$ \sourcecite{ref:d6-04-amh26}{AMH26}. Rains, Theorems~2 and~3, shows that this bound cannot be attained at length nine. Since the code dimension is an integer, it is at most $111$ \sourcecite{ref:d6-04-rains99}{Rai99}.
\end{itemize}

\subsection*{References}

\begin{enumerate}[label={},leftmargin=4.8em,labelsep=0.5em]
  \footnotesize
  \item[\textup{[AMH26]}]\label{ref:d6-04-amh26}
  G. Anglès Munné and F. Huber, ``SDP bounds on quantum codes: rational certificates,'' arXiv preprint (2026). \href{https://arxiv.org/abs/2603.19901v1}{arXiv:2603.19901v1}.

  \item[\textup{[ROJ19]}]\label{ref:d6-04-roj19}
  A. Rigby, J. C. Olivier, and P. Jarvis, ``Heuristic construction of codeword stabilized codes,'' \emph{Physical Review A} \textbf{100}, 062303 (2019). \href{https://doi.org/10.1103/PhysRevA.100.062303}{doi:10.1103/PhysRevA.100.062303}; \href{https://arxiv.org/abs/1907.04537}{arXiv:1907.04537}.

  \item[\textup{[Rai99]}]\label{ref:d6-04-rains99}
  E. M. Rains, \emph{Quantum codes of minimum distance two}, IEEE Transactions on Information Theory \textbf{45}, 266--271 (1999). \href{https://doi.org/10.1109/18.746807}{doi:10.1109/18.746807}. \href{https://arxiv.org/abs/quant-ph/9704043}{arXiv:quant-ph/9704043}.
\end{enumerate}

\subsection*{Comment}

The \href{https://qiqc-op.com/problem/op_5ea7cb9abb281829/}{seven-qubit distance-two problem} imposes the same detection requirement on a shorter block. The \href{https://qiqc-op.com/problem/op_0d6ac674da115857/}{nine-qubit distance-three problem} also requires detection of errors of weight two. The catalog and literature were checked through November 26, 2026.

\subsection*{Field}

Quantum Error Correction

\subsection*{Topic}

Quantum coding theory

\subsection*{ID}

\texttt{op\_0af3e22238a42d74}
